\section{Coin Combinations II}
\label{sec:cses-coin-combinations-ii}

\problemstatement{Given $n$ coins and a target $x$, count the number of
\emph{distinct} ways to make $x$ where the \emph{order does not matter}
(each multiset of coins is counted once). Output modulo $10^9+7$.}

\begin{patternbox}
\emph{``Count ways, order does not matter''} is the combination count.
The trick that removes duplicate orderings is to fix a global order on
the coins and only ever extend a combination with coins ``at or after''
the current one -- realised by putting the \textbf{coin loop on the
outside}.
\end{patternbox}

\subsection*{Why coin-outer counts each multiset once}

Process coins one at a time. After incorporating the first $i$ coins,
$\textit{dp}[s]$ counts the combinations of \emph{those coins only} that
sum to $s$. Adding coin $c_i$ updates
$\textit{dp}[s]\mathrel{+}=\textit{dp}[s-c_i]$ for increasing $s$, which
appends copies of $c_i$ to combinations already built from coins
$\le i$. Because a combination is built in a fixed coin order, each
multiset is generated exactly once -- never as a reordering.

\begin{ideabox}[Coin-Outer Loop Counts Combinations]
Freezing the coin order and only extending forward eliminates permutations
of the same multiset. This is the exact mirror of Coin Combinations~I:
same array update, loops swapped. Coin-outer $\Rightarrow$ combinations;
sum-outer $\Rightarrow$ permutations.
\end{ideabox}

\subsection*{Algorithm}

\begin{enumerate}
  \item $\textit{dp}[0]=1$.
  \item For each coin $c$ (outer): for $s=c\ldots x$ (inner),
        $\textit{dp}[s]\mathrel{+}=\textit{dp}[s-c]$ modulo $10^9+7$.
  \item Answer $\textit{dp}[x]$.
\end{enumerate}

\subsection*{Dry run}

\dryrun{Coins $\{2,3,5\}$, $x=9$.}

\begin{itemize}
  \item After coin $2$: $\textit{dp}[\text{even}]=1$ for $0,2,4,6,8$.
  \item After coin $3$: combinations of $\{2,3\}$; e.g.\
        $\textit{dp}[5]=1$ ($2{+}3$), $\textit{dp}[9]$ grows.
  \item After coin $5$: final $\textit{dp}[9]=3$ -- the multisets
        $\{2,2,5\}$, $\{2,2,2,3\}$, $\{3,3,3\}$.
\end{itemize}

\subsection*{C++ implementation}

\begin{lstlisting}
#include <bits/stdc++.h>
using namespace std;
const int MOD = 1e9 + 7;

int main() {
    int n, x;
    cin >> n >> x;
    vector<int> coins(n);
    for (int& c : coins) cin >> c;

    vector<long long> dp(x + 1, 0);
    dp[0] = 1;
    for (int c : coins)                   // coin outer -> combinations
        for (int s = c; s <= x; s++)
            dp[s] = (dp[s] + dp[s - c]) % MOD;

    cout << dp[x] << "\n";
    return 0;
}
\end{lstlisting}

\begin{complexitybox}
\textbf{Time Complexity.} $O(nx)$. \textbf{Space Complexity.} $O(x)$.
\end{complexitybox}

\begin{takeawaybox}
Combinations vs permutations in coin DP is decided by \emph{which loop is
outer}, nothing else. Keep the pair Coin Combinations I and II side by
side in memory as the canonical illustration -- the same one-line update,
opposite loop nesting, opposite meaning.
\end{takeawaybox}
